% Generated draft: requires mathematical and source review.
\documentclass[11pt]{article}
\usepackage[margin=25mm]{geometry}
\usepackage{fontspec}
\setmainfont{Latin Modern Roman}
\usepackage{amsmath,amssymb,mathtools,unicode-math}
\setmathfont{Latin Modern Math}
\usepackage{enumitem,needspace,xurl,hyperref}
\setlength{\parindent}{0pt}
\setlength{\parskip}{5pt}
\setlength{\emergencystretch}{3em}
\newcommand{\N}{\mathbb{N}}
\newcommand{\Z}{\mathbb{Z}}
\newcommand{\Q}{\mathbb{Q}}
\newcommand{\R}{\mathbb{R}}
\newcommand{\C}{\mathbb{C}}
\newcommand{\F}{\mathbb{F}}
\newcommand{\A}{\mathbb{A}}
\newcommand{\PP}{\mathbb{P}}
\newcommand{\E}{\mathbb{E}}
\newcommand{\eps}{\varepsilon}
\newcommand{\dd}{\,\mathrm{d}}
\begin{document}
\section*{UnsolvedMath ID: 30005165}
\textbf{Source problem code:} catalogue entry
\subsection{Definitions and mathematical statement}
The supplied conjecture concerns dispersive evolution equations in one spatial dimension with cubic defocusing nonlinearities. The record does not specify the spatial domain, the dispersion operator, the admissible nonlinearities, the meaning of defocusing, or the function space in which initial data must be small. Thus it supplies a general conjectural principle rather than a fully specified family of initial-value problems.

A schematic quantified statement is as follows. Let \(\mathcal C\) denote the intended class of equations, and, for each \(E\in\mathcal C\), let \(X_E\) denote its intended normed initial-data space. The conjecture asks whether
\[
\forall E\in\mathcal C\;\exists\varepsilon_E>0\;\forall u_0\in X_E,\qquad
\|u_0\|_{X_E}<\varepsilon_E
\ \Longrightarrow\ 
\begin{gathered}
\text{there exists a solution }u\text{ of }E\text{ with }u(0)=u_0,\\
\text{global in time and scattering in the intended sense.}
\end{gathered}
\]
Here \(\mathcal C\), \(X_E\), the solution concept, the time directions of global existence, and the scattering requirement are unresolved placeholders, not additional hypotheses supplied by the source. The formula records the universal small-data assertion without claiming a precise admissible class or a uniform threshold across equations.

For clarity, if an equation has a linear evolution group \(U_E(t)\), ordinary forward scattering in a specified normed space \(Y_E\) means that there exists \(u_+\in Y_E\) such that
\[
\lim_{t\to+\infty}\|u(t)-U_E(t)u_+\|_{Y_E}=0.
\]
Backward scattering is defined analogously using \(t\to-\infty\) and a state \(u_-\). These definitions explain one possible meaning of scattering; the supplied record does not establish that ordinary scattering, rather than modified scattering with an appropriate asymptotic correction, is intended. Neither a correction nor a scattering topology is specified.

The source background says that many global well-posedness results were obtained for particular models with initial data both small and localized. The stated conjecture assumes smallness but gives no explicit localization hypothesis. Global well-posedness in that background must not be identified with the conjectured scattering conclusion.
\subsection{Short English statement}
The conjecture asks whether every intended one-dimensional cubic defocusing dispersive equation has global scattering solutions for all sufficiently small initial data. The equation class, smallness norm, spatial domain, time directions, and meaning of scattering remain unspecified, so this is a schematic statement requiring mathematical review.
\subsection{Sources}
\begin{itemize}
\item Supplied local UnsolvedMath registry record, permanent ID \(30005165\), title ``Scattering for General Cubic Defocusing Dispersive Flows,'' source problem code \texttt{catalogue entry}. The record attributes the excerpt to Daniel Tataru's presentation ``Global solutions for 1D cubic defocusing dispersive equations,'' joint work with Mihaela Ifrim, and labels the assertion Conjecture 1. The proposed-year field is 2022; no separate proposer is named.
\item Supplied report citation: \emph{Real Analysis, Harmonic Analysis and Applications} (2023), Oberwolfach Reports, pp.\ 1743--1803, DOI \texttt{10.4171/owr/2022/31}. The literature-source label is ``Oberwolfach Report 31/2022.'' \url{https://doi.org/10.4171/owr/2022/31}. These supplied year descriptions are retained without independent bibliographic reconciliation.
\item Supplied literature source: Ifrim--Tataru, ``Global solutions, Part I.'' \url{https://doi.org/10.1017/fmp.2023.30}.
\item Supplied literature source: Ifrim--Tataru, ``Global solutions, Part IV.'' \url{https://arxiv.org/abs/2410.10052}.
\item Literature assessment dated exactly \texttt{2026-08-22}: the supplied curation labels the problem ``partially solved'' and reports that Ifrim--Tataru establish global scattering for a broad class of small, nonlocalized data, with Part IV enlarging the dispersion class, while the broad question remains unsettled for every one-dimensional cubic defocusing flow. This is supplied dated curation derived from the OWR record, not an independently verified literature assessment. The cited claims and time-sensitive status require independent verification.
\end{itemize}
\end{document}
